\subsection{Xác định tác nhân hệ thống}

Phần này xác định các tác nhân tham gia hệ thống, các ca sử dụng tổng quát tương ứng với từng tác nhân và mối quan hệ giữa các ca sử dụng trong phạm vi biểu đồ ca sử dụng. Dựa trên phạm vi triển khai, hệ thống trợ lý ảo tuyển sinh gồm bốn vai trò chính: Khách, Người dùng, Quản lý và Admin.

\subsubsection{Danh sách tác nhân}
\begin{table}[H]
\centering
\caption{Danh sách tác nhân trong hệ thống}
\label{tab:actors_list}
\begin{tabular}{|p{3.2cm}|p{6.8cm}|p{5cm}|}
\hline
\textbf{Tác nhân} & \textbf{Mô tả} & \textbf{Phạm vi sử dụng} \\ \hline
Khách & Là người truy cập hệ thống khi chưa đăng nhập, có nhu cầu hỏi đáp nhanh với trợ lý ảo & Hỏi đáp với Rasa \\ \hline
Người dùng & Là người sử dụng hệ thống sau khi đăng nhập để hỏi đáp và theo dõi thông tin hội thoại của bản thân & Hỏi đáp với Rasa, quản lý lịch sử hội thoại \\ \hline
Quản lý & Là người phụ trách vận hành dữ liệu nghiệp vụ, tài liệu, biểu mẫu, bộ dữ liệu và các cấu hình phục vụ chatbot & Quản lý bộ dữ liệu, văn bản biểu mẫu, bộ API, người dùng theo phạm vi được cấp quyền \\ \hline
Admin & Là người quản trị cấp cao, chịu trách nhiệm phân quyền, quản trị người dùng và theo dõi báo cáo thống kê toàn hệ thống & Quản lý người dùng, vai trò, quyền hạn và báo cáo thống kê \\ \hline
\end{tabular}
\end{table}

\subsubsection{Danh sách ca sử dụng tổng quát}

\begin{table}[H]
\centering
\caption{Danh sách ca sử dụng tổng quát}
\label{tab:usecase_list}
\begin{tabular}{|p{1.5cm}|p{4.2cm}|p{3cm}|p{6.3cm}|}
\hline
\textbf{Mã ca sử dụng} & \textbf{Tên ca sử dụng} & \textbf{Tác nhân chính} & \textbf{Mô tả} \\ \hline
UC-01 & Hỏi đáp với Rasa & Khách, Người dùng & Tác nhân đặt câu hỏi tuyển sinh và nhận phản hồi từ trợ lý ảo Rasa \\ \hline
UC-02 & Đăng nhập & Người dùng, Quản lý, Admin & Xác thực tài khoản trước khi sử dụng các chức năng cần phân quyền \\ \hline
UC-03 & Quản lý lịch sử hội thoại & Người dùng & Xem và quản lý các hội thoại đã thực hiện với chatbot \\ \hline
UC-04 & Quản lý các bộ dữ liệu & Quản lý & Thêm, xem chi tiết, thu thập dữ liệu và quản lý câu hỏi gợi ý \\ \hline
UC-05 & Quản lý văn bản, biểu mẫu & Quản lý & Quản lý văn bản và tài liệu ngữ cảnh phục vụ tra cứu, hỏi đáp \\ \hline
UC-06 & Quản lý bộ API & Quản lý & Quản lý các API hoặc tài liệu API phục vụ tích hợp hệ thống \\ \hline
UC-07 & Quản lý chatbot Rasa & Admin & Quản lý intent, entity, action, response, rule, story, slot và cấu hình chatbot \\ \hline
UC-08 & Quản lý người dùng & Quản lý, Admin & CRUD người dùng, ban/unban, thiết lập vai trò và cấp chatbot \\ \hline
UC-09 & Quản lý vai trò và quyền hạn & Admin & Quản lý vai trò, quyền hạn, gán quyền cho vai trò và gán vai trò cho người dùng \\ \hline
UC-10 & Báo cáo thống kê & Admin & Xem thống kê người dùng, hội thoại, chatbot, NLP và văn bản \\ \hline
\end{tabular}
\end{table}

\subsubsection{Mô tả quan hệ giữa các ca sử dụng}

Các ca sử dụng dành cho Người dùng, Quản lý và Admin đều cần thực hiện đăng nhập trước khi truy cập chức năng tương ứng. Vai trò Khách chỉ sử dụng chức năng hỏi đáp cơ bản với Rasa; khi cần lưu lịch sử hoặc dùng các chức năng cá nhân, Khách có thể chuyển thành Người dùng sau khi đăng nhập.

Nhóm chức năng của Quản lý tập trung vào dữ liệu nghiệp vụ, văn bản, biểu mẫu, bộ dữ liệu và bộ API. Nhóm chức năng của Admin tập trung vào quản trị quyền hạn, người dùng, chatbot Rasa và báo cáo thống kê. Cách phân tách này giúp hệ thống kiểm soát rõ phạm vi thao tác giữa người vận hành nội dung và người quản trị cấp cao.
